\documentclass[aspectratio=169]{beamer}
\usetheme{Madrid}
\usecolortheme{default}
\usepackage[utf8]{inputenc}
\usepackage{amsmath}
\usepackage{amsfonts}
\usepackage{amssymb}
\usepackage{enumerate}

\setbeamertemplate{navigation symbols}{}
\setbeamerfont{block body}{size=\small}
\setbeamerfont{alertblock body}{size=\small}

\title{CRISC Governance, Ethics and Analysis -- 100 Real-Life Questions}
\subtitle{Governance, Ethics, Threats, Analysis and Exam Mindset -- Practical Scenarios}
\author{CRISC Mastery Series}
\date{\today}

% ------------------------------------------------------------
% \hq{number}{topic}{question}{options}{letter}{explanation}{takeaway}
% Creates one question frame and one explanation frame.
% ------------------------------------------------------------
\newcommand{\hq}[7]{%
	\begin{frame}{Q#1: #2}
		\begin{block}{Question}
			#3
		\end{block}
		\begin{itemize}
			#4
		\end{itemize}
	\end{frame}
	\begin{frame}{Q#1: Explanation}
		\begin{block}{Correct Answer: #5}
			\begin{itemize}
				#6
			\end{itemize}
		\end{block}
		\begin{alertblock}{Key Takeaway}
			#7
		\end{alertblock}
	\end{frame}
}

\begin{document}
	
	\begin{frame}
		\titlepage
	\end{frame}
	
	\begin{frame}{Instructions}
		\begin{alertblock}{How to Use This Deck}
			Think like a risk practitioner in a real organization.\\
			Focus on business context, ownership, and practical decision-making.\\
			The answer and a key takeaway are on the slide following each question.
		\end{alertblock}
	\end{frame}
	
	\begin{frame}{Outline}
		\tableofcontents
	\end{frame}
	
	% ============================================================
	\section{Governance, Strategy and Culture}
	% ============================================================
	
	\hq{1}{Governance vs Management}
	{Which statement \textbf{BEST} distinguishes governance from management?}
	{\item[\textbf{(A)}] They are the same function with different names
		\item[\textbf{(B)}] Governance runs daily operations, and management sets strategy
		\item[\textbf{(C)}] Governance sets direction and oversight, while management plans, builds, and runs activities to achieve it
		\item[\textbf{(D)}] Governance applies only to IT, and management applies to finance}
	{C}
	{\item The board and executives direct and oversee through governance.
		\item Management carries out the direction day to day.
		\item Confusing the two blurs accountability.}
	{Governance directs; management executes.}
	
	\hq{2}{Aligning Risk to Strategy}
	{The company adopts a strategy of rapid digital expansion. What should the risk practitioner do \textbf{FIRST}?}
	{\item[\textbf{(A)}] Freeze all new technology projects
		\item[\textbf{(B)}] Copy the risk register of a competitor
		\item[\textbf{(C)}] Understand the strategic objectives and identify the risks that could affect them
		\item[\textbf{(D)}] Increase the security budget by a fixed amount}
	{C}
	{\item Risk is defined relative to objectives.
		\item Understanding strategy lets you focus on what matters.
		\item Blanket restrictions can block the strategy itself.}
	{Risk management supports objectives; it does not oppose them.}
	
	\hq{3}{Board Oversight}
	{What is the board's \textbf{PRIMARY} role in IT risk?}
	{\item[\textbf{(A)}] Configure firewalls
		\item[\textbf{(B)}] Approve every access request
		\item[\textbf{(C)}] Perform vulnerability scans
		\item[\textbf{(D)}] Provide oversight and ensure risk management aligns with strategy and appetite}
	{D}
	{\item The board is accountable for governance.
		\item It sets or endorses risk appetite.
		\item Operational tasks belong to management and specialists.}
	{The board oversees; it does not operate.}
	
	\hq{4}{Risk Culture}
	{Employees hide mistakes because they fear punishment. What is the \textbf{BEST} way to improve the risk culture?}
	{\item[\textbf{(A)}] Add stricter penalties for every error
		\item[\textbf{(B)}] Publish a longer policy
		\item[\textbf{(C)}] Leaders model transparency and reward timely reporting of issues
		\item[\textbf{(D)}] Limit reporting to managers only}
	{C}
	{\item Culture follows what leaders do and reward.
		\item Fear suppresses reporting and increases hidden risk.
		\item Policies alone do not change behavior.}
	{People report risk where it is safe to do so.}
	
	\hq{5}{Policy Hierarchy}
	{Which describes the correct relationship among policies, standards, and procedures?}
	{\item[\textbf{(A)}] Procedures set strategy, and policies list steps
		\item[\textbf{(B)}] Standards replace policies
		\item[\textbf{(C)}] They all contain the same content
		\item[\textbf{(D)}] Policies state intent, standards set mandatory requirements, and procedures give step-by-step instructions}
	{D}
	{\item Policies are high-level and approved by management.
		\item Standards specify measurable requirements.
		\item Procedures explain how to carry out tasks.}
	{Why, what, then how.}
	
	\hq{6}{Outdated Policy}
	{A security policy has not been reviewed in six years and no longer matches the technology in use. What is the \textbf{BEST} action?}
	{\item[\textbf{(A)}] Review and update it through the defined approval process, and communicate the changes
		\item[\textbf{(B)}] Leave it because it was once approved
		\item[\textbf{(C)}] Delete it
		\item[\textbf{(D)}] Let each team write its own}
	{A}
	{\item Policies must stay relevant to the environment.
		\item Approval gives them authority.
		\item Communication ensures awareness.}
	{A stale policy is a weak control.}
	
	\hq{7}{RACI Gap}
	{During an incident, two teams each assume the other is responsible for communication. Which tool would have \textbf{BEST} prevented this?}
	{\item[\textbf{(A)}] A larger incident budget
		\item[\textbf{(B)}] A RACI matrix defining who is responsible, accountable, consulted, and informed
		\item[\textbf{(C)}] A new firewall
		\item[\textbf{(D)}] A longer password policy}
	{B}
	{\item RACI clarifies roles for each activity.
		\item Exactly one party should be accountable.
		\item Clear roles speed response.}
	{If everyone owns it, no one does.}
	
	\hq{8}{Asset Inventory First}
	{A risk assessment is starting, but the asset inventory is incomplete. What is the \textbf{BEST} next step?}
	{\item[\textbf{(A)}] Assess only the assets already known to IT
		\item[\textbf{(B)}] Skip assets and list threats only
		\item[\textbf{(C)}] Identify and prioritize critical assets with their owners before assessing risk
		\item[\textbf{(D)}] Assess every device at equal depth}
	{C}
	{\item You cannot assess risk to assets you do not know.
		\item Prioritization focuses effort on business value.
		\item Owners provide value and impact information.}
	{No inventory, no meaningful assessment.}
	
	\hq{9}{Asset Valuation}
	{Which factor is \textbf{MOST} important when valuing an information asset?}
	{\item[\textbf{(A)}] Its importance to business processes and the impact if it were lost, altered, or disclosed
		\item[\textbf{(B)}] The purchase price of the hardware
		\item[\textbf{(C)}] The age of the server
		\item[\textbf{(D)}] The brand of the software}
	{A}
	{\item Information value depends on business use.
		\item Replacement cost can be much lower than impact.
		\item Value drives the level of protection.}
	{Value the impact, not the invoice.}
	
	\hq{10}{Risk Strategy}
	{Which of these \textbf{BEST} describes an IT risk strategy?}
	{\item[\textbf{(A)}] A list of all software licenses
		\item[\textbf{(B)}] A plan, aligned with business objectives, for how the organization will identify, assess, respond to, and monitor IT risk
		\item[\textbf{(C)}] A schedule of penetration tests
		\item[\textbf{(D)}] A report of last year's incidents}
	{B}
	{\item Strategy sets the overall approach to risk.
		\item It supports business objectives.
		\item Lists and reports are inputs, not strategy.}
	{Strategy answers how risk will be managed, and why.}
	
	\hq{11}{IT Governance Frameworks}
	{The organization wants a framework to align IT with business goals and manage IT governance. Which is the \textbf{BEST} fit?}
	{\item[\textbf{(A)}] A network cabling standard
		\item[\textbf{(B)}] A project scheduling tool
		\item[\textbf{(C)}] A spreadsheet template
		\item[\textbf{(D)}] COBIT}
	{D}
	{\item COBIT focuses on governance and management of enterprise IT.
		\item It links business goals to IT goals.
		\item Frameworks should be adapted to the organization.}
	{COBIT governs IT; do not confuse it with tools.}
	
	\hq{12}{Cybersecurity Framework Use}
	{A company wants a common language for describing its cybersecurity activities from identification to recovery. Which framework is \textbf{MOST} suitable?}
	{\item[\textbf{(A)}] A software license agreement
		\item[\textbf{(B)}] A coding style guide
		\item[\textbf{(C)}] A procurement catalog
		\item[\textbf{(D)}] NIST Cybersecurity Framework}
	{D}
	{\item The NIST CSF organizes activities into functions.
		\item It helps communicate across technical and business teams.
		\item It supports maturity assessment.}
	{Use a framework to speak one language.}
	
	\hq{13}{IT Service Management}
	{Which framework is \textbf{MOST} associated with IT service management practices such as incident and change management?}
	{\item[\textbf{(A)}] ITIL
		\item[\textbf{(B)}] The CRISC code of ethics
		\item[\textbf{(C)}] A firewall rule base
		\item[\textbf{(D)}] A data classification scheme}
	{A}
	{\item ITIL describes service management processes.
		\item COBIT governs; ITIL manages services.
		\item They can be used together.}
	{COBIT governs, ITIL runs the services.}
	
	\hq{14}{Risk Practitioner Role}
	{What is the risk practitioner's \textbf{PRIMARY} role in governance?}
	{\item[\textbf{(A)}] Make all risk decisions for the business
		\item[\textbf{(B)}] Provide independent, objective risk information so decision makers can make informed choices
		\item[\textbf{(C)}] Approve all projects
		\item[\textbf{(D)}] Operate the controls}
	{B}
	{\item Owners make decisions; practitioners inform them.
		\item Practitioners help identify, assess, and report risk.
		\item Operating controls creates conflicts of interest.}
	{Advise and inform; do not decide for the owner.}
	
	\hq{15}{Competing Priorities}
	{Two departments demand opposite risk responses for a shared system. What should the risk practitioner do?}
	{\item[\textbf{(A)}] Pick the department that complained first
		\item[\textbf{(B)}] Implement both responses at once
		\item[\textbf{(C)}] Delay the decision indefinitely
		\item[\textbf{(D)}] Present the risk, business impact, and options to the accountable executive or governance body for decision}
	{D}
	{\item Conflicts need an accountable decision maker.
		\item Clear impact information supports a fair choice.
		\item Delay leaves the risk open.}
	{Escalate conflicts with facts, not opinions.}
	
	\hq{16}{Ethical Culture Link}
	{How are ethics and risk culture related?}
	{\item[\textbf{(A)}] They are unrelated
		\item[\textbf{(B)}] Ethics matters only for external auditors
		\item[\textbf{(C)}] Culture changes only through technology
		\item[\textbf{(D)}] An ethical culture encourages honest reporting, whereas an unethical one hides or manipulates risk information}
	{D}
	{\item Behavior at the top influences behavior throughout.
		\item Honest reporting depends on trust.
		\item Ethical lapses can create hidden risk.}
	{Ethics shapes what gets reported.}
	
	% ============================================================
	\section{Three Lines, Appetite and Roles}
	% ============================================================
	
	\hq{17}{First Line of Defense}
	{Who forms the first line of defense in the three lines model?}
	{\item[\textbf{(A)}] Internal audit
		\item[\textbf{(B)}] Business and operational management who own and manage risks and controls
		\item[\textbf{(C)}] The board's audit committee
		\item[\textbf{(D)}] External regulators}
	{B}
	{\item The first line owns risk in daily operations.
		\item The second line oversees and supports.
		\item The third line provides independent assurance.}
	{First line owns the risk.}
	
	\hq{18}{Second Line}
	{Which function is a typical second line of defense?}
	{\item[\textbf{(A)}] Internal audit
		\item[\textbf{(B)}] Risk management and compliance functions that oversee and challenge the first line
		\item[\textbf{(C)}] The sales team
		\item[\textbf{(D)}] External auditors}
	{B}
	{\item The second line sets frameworks and monitors.
		\item It advises but does not own the operations.
		\item Independent assurance belongs to the third line.}
	{Second line oversees and challenges.}
	
	\hq{19}{Third Line}
	{What is the role of internal audit in the three lines model?}
	{\item[\textbf{(A)}] Operate preventive controls
		\item[\textbf{(B)}] Own and accept business risks
		\item[\textbf{(C)}] Design the risk appetite
		\item[\textbf{(D)}] Provide independent assurance on governance, risk management, and control effectiveness to the board}
	{D}
	{\item Independence is essential to assurance.
		\item Audit reports to the board or audit committee.
		\item Operating controls would compromise objectivity.}
	{Third line assures; it does not operate.}
	
	\hq{20}{Independence Concern}
	{Internal audit is asked to implement a new control and then audit it. What is the \textbf{MAIN} concern?}
	{\item[\textbf{(A)}] Higher cost
		\item[\textbf{(B)}] Slower deployment
		\item[\textbf{(C)}] Fewer tools
		\item[\textbf{(D)}] Loss of independence and objectivity}
	{D}
	{\item Auditing your own work creates a conflict.
		\item Independence supports credible assurance.
		\item Management should implement controls.}
	{Do not audit what you built.}
	
	\hq{21}{Risk Appetite Basics}
	{Which statement about risk appetite is \textbf{MOST} accurate?}
	{\item[\textbf{(A)}] It is the amount and type of risk the organization is willing to pursue to achieve its objectives
		\item[\textbf{(B)}] It is the maximum risk the organization can survive
		\item[\textbf{(C)}] It is the risk remaining after controls
		\item[\textbf{(D)}] It is the cost of insurance}
	{A}
	{\item Appetite is a strategic choice set by leadership.
		\item Capacity is the maximum survivable level.
		\item Residual risk is measured after controls.}
	{Appetite is willingness; capacity is ability.}
	
	\hq{22}{Tolerance vs Appetite}
	{Which describes risk tolerance?}
	{\item[\textbf{(A)}] The acceptable deviation from appetite for specific objectives, often expressed as thresholds
		\item[\textbf{(B)}] The amount of insurance purchased
		\item[\textbf{(C)}] The number of controls in place
		\item[\textbf{(D)}] The average vulnerability age}
	{A}
	{\item Tolerance translates appetite into measurable limits.
		\item It is specific to objectives or risk types.
		\item KRIs often track tolerance thresholds.}
	{Tolerance puts numbers on appetite.}
	
	\hq{23}{Risk Capacity}
	{What does risk capacity represent?}
	{\item[\textbf{(A)}] The risk that management prefers
		\item[\textbf{(B)}] The risk after mitigation
		\item[\textbf{(C)}] The risk transferred to insurers
		\item[\textbf{(D)}] The maximum risk the organization can bear before its viability is threatened}
	{D}
	{\item Capacity is a limit, not a preference.
		\item Appetite should be set below capacity.
		\item Exceeding capacity threatens survival.}
	{Never set appetite above capacity.}
	
	\hq{24}{Risk Profile}
	{What is a risk profile?}
	{\item[\textbf{(A)}] A single risk's remediation plan
		\item[\textbf{(B)}] A description of the organization's identified risks and their levels at a point in time
		\item[\textbf{(C)}] A staff performance review
		\item[\textbf{(D)}] A vendor's marketing profile}
	{B}
	{\item The profile summarizes risk exposure.
		\item It is compared against appetite and tolerance.
		\item It changes as risks and controls change.}
	{The risk profile is the current picture of exposure.}
	
	\hq{25}{Acceptance Documentation}
	{A risk owner verbally agrees to accept a high risk. What should the risk practitioner do?}
	{\item[\textbf{(A)}] Rely on the verbal agreement
		\item[\textbf{(B)}] Ask audit to sign instead
		\item[\textbf{(C)}] Ensure the acceptance is formally documented and approved at the appropriate authority level
		\item[\textbf{(D)}] Mark the risk closed}
	{C}
	{\item Accepted risk needs a record for accountability.
		\item Authority levels should match the risk level.
		\item Closed means resolved, not accepted.}
	{If it is not written down, it was not accepted.}
	
	\hq{26}{Accepting Beyond Authority}
	{A department manager wants to accept a risk that exceeds their authority limit. What should happen?}
	{\item[\textbf{(A)}] Let the manager accept it anyway
		\item[\textbf{(B)}] Reduce the rating to fit the limit
		\item[\textbf{(C)}] Escalate to the level of management that has authority for that level of risk
		\item[\textbf{(D)}] Ignore the request}
	{C}
	{\item Authority limits reflect governance decisions.
		\item Adjusting ratings to fit limits is manipulation.
		\item Escalation maintains accountability.}
	{Risk acceptance must go to the right level.}
	
	\hq{27}{Risk Practitioner vs Owner}
	{Who is typically accountable for the day-to-day management of a specific risk?}
	{\item[\textbf{(A)}] The risk practitioner
		\item[\textbf{(B)}] The external auditor
		\item[\textbf{(C)}] The vendor
		\item[\textbf{(D)}] The risk owner}
	{D}
	{\item Owners have authority over the related process and resources.
		\item Practitioners provide guidance and challenge.
		\item Accountability cannot be delegated outward.}
	{Risk owner = accountable.}
	
	\hq{28}{Control Owner vs Risk Owner}
	{Which statement correctly differentiates a control owner from a risk owner?}
	{\item[\textbf{(A)}] The control owner operates a control, and the risk owner is accountable for the risk the control addresses
		\item[\textbf{(B)}] They are always the same person
		\item[\textbf{(C)}] The control owner accepts all risk
		\item[\textbf{(D)}] The risk owner operates all controls}
	{A}
	{\item Control owners perform and evidence controls.
		\item Risk owners decide on response and acceptance.
		\item Roles can overlap but responsibilities differ.}
	{One runs the control, one owns the risk.}
	
	\hq{29}{Appetite Cascades}
	{Different departments define their own conflicting risk appetites. What is the \textbf{BEST} approach?}
	{\item[\textbf{(A)}] Let each department set its own without review
		\item[\textbf{(B)}] Use the most conservative one only
		\item[\textbf{(C)}] Remove appetite statements
		\item[\textbf{(D)}] Align them to an enterprise-level appetite approved by leadership and cascade it consistently}
	{D}
	{\item Enterprise appetite ensures consistency.
		\item Departments set tolerances within it.
		\item Conflicts weaken decisions and reporting.}
	{One enterprise appetite, many aligned tolerances.}
	
	\hq{30}{Escalating Cyber Risk}
	{Which cybersecurity issue is \textbf{MOST} appropriate for escalation to the board?}
	{\item[\textbf{(A)}] A risk exceeding the approved appetite that threatens strategic objectives
		\item[\textbf{(B)}] A single low-severity virus alert
		\item[\textbf{(C)}] A routine patch schedule
		\item[\textbf{(D)}] A minor printer outage}
	{A}
	{\item Boards need strategic, material information.
		\item Low-level events are handled operationally.
		\item Escalation criteria should be defined.}
	{Escalate what is material, not what is merely noisy.}
	
	% ============================================================
	\section{Laws, Ethics and Professional Conduct}
	% ============================================================
	
	\hq{31}{Applicable Regulations}
	{The company expands into a new region. What should the risk practitioner do \textbf{FIRST} regarding legal requirements?}
	{\item[\textbf{(A)}] Assume home-country laws apply everywhere
		\item[\textbf{(B)}] Work with legal and compliance to identify applicable laws and their impact on processes and controls
		\item[\textbf{(C)}] Wait for a regulator's letter
		\item[\textbf{(D)}] Rely on a competitor's website}
	{B}
	{\item Legal obligations vary by jurisdiction.
		\item Legal experts interpret the requirements.
		\item Early identification avoids costly rework.}
	{Know which laws apply before you operate.}
	
	\hq{32}{Compliance vs Risk}
	{A control meets a regulation, but a risk assessment still shows high exposure. What does this indicate?}
	{\item[\textbf{(A)}] Compliance always means no risk
		\item[\textbf{(B)}] Compliance does not guarantee acceptable risk, so additional response may be needed
		\item[\textbf{(C)}] The risk assessment is wrong
		\item[\textbf{(D)}] The regulation is unnecessary}
	{B}
	{\item Regulations set minimum requirements.
		\item Risk depends on the specific context.
		\item Both views are needed.}
	{Compliant is not the same as secure.}
	
	\hq{33}{Code of Ethics}
	{A risk practitioner discovers that management is overstating control effectiveness in a report to the board. What should they do?}
	{\item[\textbf{(A)}] Stay silent to protect the manager
		\item[\textbf{(B)}] Leak the report externally at once
		\item[\textbf{(C)}] Raise the concern through appropriate channels and ensure accurate information is reported, in line with professional ethics
		\item[\textbf{(D)}] Change the report secretly to match what they believe}
	{C}
	{\item Professional codes require honesty and integrity.
		\item Use proper escalation channels first.
		\item Unilateral or secret action can create further problems.}
	{Integrity first; escalate through proper channels.}
	
	\hq{34}{Conflict of Interest}
	{A risk practitioner is asked to assess a vendor owned by a close relative. What is the \textbf{BEST} action?}
	{\item[\textbf{(A)}] Do the assessment and say nothing
		\item[\textbf{(B)}] Disclose the conflict and recuse themselves from the assessment
		\item[\textbf{(C)}] Give the vendor a favorable rating
		\item[\textbf{(D)}] Delay the project}
	{B}
	{\item Undisclosed conflicts damage trust.
		\item Independence supports objective assessment.
		\item Disclosure lets management reassign the task.}
	{Disclose conflicts early.}
	
	\hq{35}{Confidentiality Duty}
	{A practitioner learns confidential information during an assessment. When may it be shared?}
	{\item[\textbf{(A)}] With anyone interested in the topic
		\item[\textbf{(B)}] On social media in general terms
		\item[\textbf{(C)}] With a future employer
		\item[\textbf{(D)}] Only with authorized parties, and as required by law or approved by the organization}
	{D}
	{\item Professional ethics require protecting confidential information.
		\item Disclosure must be authorized or legally required.
		\item Sharing for personal gain is a violation.}
	{What you learn on the job stays on the job.}
	
	\hq{36}{Data Breach Law}
	{A breach involving personal data occurs. What should the organization do regarding legal obligations?}
	{\item[\textbf{(A)}] Wait until the investigation is fully complete
		\item[\textbf{(B)}] Notify only if customers complain
		\item[\textbf{(C)}] Avoid documenting the breach
		\item[\textbf{(D)}] Determine notification duties and timelines with legal counsel and act within the required timeframe}
	{D}
	{\item Laws may set strict notification deadlines.
		\item Legal counsel interprets requirements.
		\item Documentation supports compliance.}
	{Notification clocks start early; plan for them.}
	
	\hq{37}{Payment Card Standard}
	{The company processes card payments. Which requirement set is \textbf{MOST} directly relevant?}
	{\item[\textbf{(A)}] A software coding style guide
		\item[\textbf{(B)}] The corporate dress code
		\item[\textbf{(C)}] The Payment Card Industry Data Security Standard
		\item[\textbf{(D)}] An office safety manual}
	{C}
	{\item PCI DSS sets security requirements for card data.
		\item Compliance is often contractual.
		\item Scope reduction lowers cost and risk.}
	{Card data brings card rules.}
	
	\hq{38}{Health Data Rules}
	{A software provider will handle patient health records for US healthcare organizations. Which law is \textbf{MOST} relevant?}
	{\item[\textbf{(A)}] The Health Insurance Portability and Accountability Act (HIPAA)
		\item[\textbf{(B)}] The Sarbanes-Oxley Act only
		\item[\textbf{(C)}] A local building code
		\item[\textbf{(D)}] The Fair Credit Reporting Act only}
	{A}
	{\item HIPAA protects health information.
		\item Business associates have obligations too.
		\item Contracts should define responsibilities.}
	{Match the data type to its governing law.}
	
	\hq{39}{Financial Reporting Controls}
	{A public company must show effective controls over financial reporting. Which law is \textbf{MOST} associated with this requirement?}
	{\item[\textbf{(A)}] The Health Insurance Portability and Accountability Act
		\item[\textbf{(B)}] The Sarbanes-Oxley Act
		\item[\textbf{(C)}] The Children's Online Privacy Protection Act
		\item[\textbf{(D)}] The Telephone Consumer Protection Act}
	{B}
	{\item SOX focuses on financial reporting and internal controls.
		\item IT general controls support those reports.
		\item Executives certify the controls.}
	{SOX equals financial reporting controls.}
	
	\hq{40}{Privacy Regulation Scope}
	{A non-EU company offers services to individuals in the EU and processes their personal data. What is true?}
	{\item[\textbf{(A)}] It never applies outside the EU
		\item[\textbf{(B)}] It applies only to EU-based companies
		\item[\textbf{(C)}] It applies only to government agencies
		\item[\textbf{(D)}] The EU General Data Protection Regulation may apply because of the targeting of EU residents}
	{D}
	{\item GDPR has extraterritorial reach.
		\item Targeting or monitoring EU individuals triggers it.
		\item Penalties can be significant.}
	{Where your customers are matters as much as where you are.}
	
	\hq{41}{Speaking Up}
	{A colleague asks the practitioner to downgrade a risk rating because the project will otherwise be delayed. What should the practitioner do?}
	{\item[\textbf{(A)}] Lower the rating to help the project
		\item[\textbf{(B)}] Decline, explain the basis for the rating, and escalate if pressure continues
		\item[\textbf{(C)}] Delete the risk from the register
		\item[\textbf{(D)}] Ignore the colleague silently}
	{B}
	{\item Ratings must reflect evidence, not convenience.
		\item Transparency protects the organization.
		\item Escalation handles undue pressure.}
	{Ratings follow facts, not deadlines.}
	
	\hq{42}{Professional Competence}
	{A practitioner is asked to assess a technology they do not understand. What is the \textbf{BEST} approach?}
	{\item[\textbf{(A)}] Give a confident opinion anyway
		\item[\textbf{(B)}] Refuse all future assessments
		\item[\textbf{(C)}] Acknowledge the limitation and obtain expert support or training before giving conclusions
		\item[\textbf{(D)}] Copy another company's rating}
	{C}
	{\item Ethics require competence and honesty.
		\item Expert input improves accuracy.
		\item Overstating knowledge creates risk.}
	{Know what you do not know.}
	
	% ============================================================
	\section{Threats, Vulnerabilities and Analysis}
	% ============================================================
	
	\hq{43}{Threat vs Vulnerability}
	{Which statement correctly relates threat, vulnerability, and risk?}
	{\item[\textbf{(A)}] A vulnerability is the same as a threat
		\item[\textbf{(B)}] Risk arises when a threat can exploit a vulnerability to cause impact on an asset
		\item[\textbf{(C)}] Risk exists without any vulnerability or threat
		\item[\textbf{(D)}] Threats are always internal}
	{B}
	{\item Threats are potential causes of harm.
		\item Vulnerabilities are weaknesses that can be exploited.
		\item Risk combines likelihood and impact.}
	{No exploitable weakness and no threat means no risk.}
	
	\hq{44}{Threat Modeling Purpose}
	{What is the primary purpose of threat modeling?}
	{\item[\textbf{(A)}] Identify how a system could be attacked and prioritize defenses early in design
		\item[\textbf{(B)}] Replace vulnerability scanning
		\item[\textbf{(C)}] Document staff vacation plans
		\item[\textbf{(D)}] Measure network bandwidth}
	{A}
	{\item It structures thinking about attackers and assets.
		\item Early use improves design.
		\item It complements, not replaces, testing.}
	{Think like the attacker before it ships.}
	
	\hq{45}{STRIDE}
	{A team categorizes threats as spoofing, tampering, repudiation, information disclosure, denial of service, and elevation of privilege. Which method is this?}
	{\item[\textbf{(A)}] SWOT
		\item[\textbf{(B)}] PERT
		\item[\textbf{(C)}] RACI
		\item[\textbf{(D)}] STRIDE}
	{D}
	{\item STRIDE is a threat categorization model.
		\item It helps systematically find threats.
		\item Other methods include attack trees and PASTA.}
	{STRIDE lists the six threat types.}
	
	\hq{46}{Vulnerability Scanning Limits}
	{A scan reports no critical vulnerabilities. What is the \textbf{MOST} accurate conclusion?}
	{\item[\textbf{(A)}] The system is completely secure
		\item[\textbf{(B)}] No attacker can succeed
		\item[\textbf{(C)}] Known vulnerabilities within the scan scope were not found, but other weaknesses may still exist
		\item[\textbf{(D)}] No more testing is ever needed}
	{C}
	{\item Scans detect known issues in scope.
		\item Logic flaws and misconfigurations can be missed.
		\item Combine scans with other assessments.}
	{A clean scan is not a clean bill of health.}
	
	\hq{47}{Penetration Test Value}
	{Why conduct a penetration test in addition to vulnerability scanning?}
	{\item[\textbf{(A)}] To avoid fixing vulnerabilities
		\item[\textbf{(B)}] To replace risk assessments
		\item[\textbf{(C)}] To test employee punctuality
		\item[\textbf{(D)}] To attempt to exploit weaknesses and show realistic attack impact}
	{D}
	{\item Pen tests demonstrate real exploitability.
		\item They can chain issues like an attacker.
		\item Results still need risk-based prioritization.}
	{Scanning finds; testing proves.}
	
	\hq{48}{Patch Prioritization}
	{Hundreds of vulnerabilities are open and patching capacity is limited. How should they be prioritized?}
	{\item[\textbf{(A)}] By severity, exploitability, asset criticality, and exposure
		\item[\textbf{(B)}] By alphabetical order of the system name
		\item[\textbf{(C)}] By the order discovered
		\item[\textbf{(D)}] By how easy they are to describe}
	{A}
	{\item Risk-based prioritization focuses on real exposure.
		\item Active exploitation raises urgency.
		\item Critical and exposed assets come first.}
	{Patch by risk, not by list order.}
	
	\hq{49}{Zero-Day}
	{A newly disclosed vulnerability has no patch yet and is being actively exploited. What is the \textbf{BEST} immediate response?}
	{\item[\textbf{(A)}] Wait quietly for a patch
		\item[\textbf{(B)}] Apply compensating controls such as isolation, filtering, or monitoring, and track vendor guidance
		\item[\textbf{(C)}] Ignore it until the next audit
		\item[\textbf{(D)}] Shut down the whole company permanently}
	{B}
	{\item Compensating controls reduce exposure.
		\item Monitoring detects exploitation attempts.
		\item A decision on temporary shutdown depends on impact.}
	{No patch? Compensate.}
	
	\hq{50}{Insider Threat}
	{Which control \textbf{BEST} reduces the risk of malicious insider data theft?}
	{\item[\textbf{(A)}] Longer lunch breaks
		\item[\textbf{(B)}] Posters about loyalty only
		\item[\textbf{(C)}] Removing all logging
		\item[\textbf{(D)}] Least privilege, activity monitoring, and data loss prevention}
	{D}
	{\item Limiting access limits possible damage.
		\item Monitoring deters and detects misuse.
		\item Technical and procedural controls work together.}
	{Trust, but limit and monitor.}
	
	\hq{51}{Social Engineering}
	{An attacker phones the help desk pretending to be an executive and asks for a password reset. What is the \textbf{BEST} defense?}
	{\item[\textbf{(A)}] More firewalls
		\item[\textbf{(B)}] Strict identity verification procedures and staff training
		\item[\textbf{(C)}] Longer passwords only
		\item[\textbf{(D)}] Faster ticket closure}
	{B}
	{\item The attack targets people, not technology.
		\item Verification steps prevent impersonation.
		\item Training builds awareness.}
	{Humans are part of the control environment.}
	
	\hq{52}{Supply Chain Attack}
	{A trusted software vendor's update contains malicious code. Which control \textbf{BEST} reduces this risk?}
	{\item[\textbf{(A)}] Install updates on all systems instantly without review
		\item[\textbf{(B)}] Stop applying all updates
		\item[\textbf{(C)}] Validate update integrity, apply staged deployment, and monitor behavior after updates
		\item[\textbf{(D)}] Trust all signed updates unconditionally}
	{C}
	{\item Supply chain attacks abuse trust.
		\item Staging limits blast radius.
		\item Monitoring detects unexpected behavior.}
	{Trust, then verify, even for vendors.}
	
	\hq{53}{Ransomware Resilience}
	{Which measure is \textbf{MOST} effective at reducing the impact of ransomware?}
	{\item[\textbf{(A)}] Tested, offline or immutable backups combined with prevention and detection controls
		\item[\textbf{(B)}] Paying the ransom as the main plan
		\item[\textbf{(C)}] Storing backups on the same network share
		\item[\textbf{(D)}] Relying on cyber insurance alone}
	{A}
	{\item Backups enable recovery without paying.
		\item Offline or immutable copies resist encryption.
		\item Prevention reduces the chance of infection.}
	{Recoverability beats negotiation.}
	
	\hq{54}{Denial of Service}
	{An online store fears downtime from large traffic floods. Which control is \textbf{MOST} relevant?}
	{\item[\textbf{(A)}] Stronger passwords
		\item[\textbf{(B)}] Encrypting backups
		\item[\textbf{(C)}] DDoS protection services and capacity planning
		\item[\textbf{(D)}] Disk quotas}
	{C}
	{\item DoS threatens availability.
		\item Mitigation services absorb or filter traffic.
		\item Response plans define actions.}
	{Match the control to the threat's target.}
	
	\hq{55}{Risk Scenarios}
	{Why are risk scenarios useful in risk assessment?}
	{\item[\textbf{(A)}] They replace the need for data
		\item[\textbf{(B)}] They describe how a threat could exploit a weakness to cause business impact, making risk understandable
		\item[\textbf{(C)}] They are used only for audits
		\item[\textbf{(D)}] They eliminate the need for owners}
	{B}
	{\item Scenarios link cause, event, and consequence.
		\item They help stakeholders discuss realistic risks.
		\item They support likelihood and impact estimates.}
	{Tell the risk as a story with consequences.}
	
	\hq{56}{Threat Intelligence}
	{How does threat intelligence \textbf{BEST} support risk assessment?}
	{\item[\textbf{(A)}] By replacing internal monitoring
		\item[\textbf{(B)}] By guaranteeing prevention
		\item[\textbf{(C)}] By providing current information on actors, techniques, and trends relevant to the organization
		\item[\textbf{(D)}] By lowering all risk ratings}
	{C}
	{\item It informs likelihood estimates.
		\item Relevance to the sector and technology matters.
		\item It supports proactive control updates.}
	{Intelligence informs likelihood.}
	
	\hq{57}{Attack Surface}
	{What is the \textbf{BEST} way to reduce an organization's attack surface?}
	{\item[\textbf{(A)}] Add more internet-facing services
		\item[\textbf{(B)}] Remove unnecessary systems, services, and access that are not needed for the business
		\item[\textbf{(C)}] Keep old accounts active
		\item[\textbf{(D)}] Disable patching}
	{B}
	{\item Fewer exposed components mean fewer opportunities.
		\item Unused assets add risk without benefit.
		\item Inventory is needed to find them.}
	{What is not there cannot be attacked.}
	
	\hq{58}{Vulnerability Management Program}
	{Which element is essential to a mature vulnerability management program?}
	{\item[\textbf{(A)}] A single annual scan
		\item[\textbf{(B)}] A repeatable cycle of discovery, prioritization, remediation, verification, and reporting
		\item[\textbf{(C)}] A list with no owners
		\item[\textbf{(D)}] Fixing only what auditors mention}
	{B}
	{\item Vulnerabilities appear continually.
		\item Verification confirms the fix worked.
		\item Reporting supports accountability.}
	{Vulnerability management is a cycle, not an event.}
	
	% ============================================================
	\section{Assessment Methods, BIA and Analysis}
	% ============================================================
	
	\hq{59}{Qualitative vs Quantitative}
	{The organization has little historical loss data and needs a quick, consistent view of many risks. Which approach is \textbf{MOST} appropriate?}
	{\item[\textbf{(A)}] Quantitative modeling with precise dollar values for all risks
		\item[\textbf{(B)}] No assessment until data exists
		\item[\textbf{(C)}] Qualitative assessment using defined rating scales
		\item[\textbf{(D)}] Random scoring}
	{C}
	{\item Qualitative methods work with limited data.
		\item Defined scales improve consistency.
		\item Quantitative analysis may follow for key risks.}
	{Use the method the available data supports.}
	
	\hq{60}{Quantitative Strength}
	{When is a quantitative risk assessment \textbf{MOST} valuable?}
	{\item[\textbf{(A)}] When financial justification is needed to compare options for high-impact risks and reliable data is available
		\item[\textbf{(B)}] When there is no data at all
		\item[\textbf{(C)}] For every low-level risk
		\item[\textbf{(D)}] When the audience dislikes numbers}
	{A}
	{\item Dollar values support cost-benefit decisions.
		\item Quality of results depends on data quality.
		\item Use for material risks where effort is justified.}
	{Numbers help when you can trust the inputs.}
	
	\hq{61}{Annualized Loss}
	{A single loss event is estimated at 200,000 dollars and is expected to occur once every four years. What is the annualized loss expectancy?}
	{\item[\textbf{(A)}] 50,000 dollars
		\item[\textbf{(B)}] 200,000 dollars
		\item[\textbf{(C)}] 800,000 dollars
		\item[\textbf{(D)}] 4,000 dollars}
	{A}
	{\item ALE equals single loss expectancy times annual rate of occurrence.
		\item Once every four years means a rate of 0.25 per year.
		\item 200,000 multiplied by 0.25 equals 50,000.}
	{ALE = SLE times ARO.}
	
	\hq{62}{Control Cost Decision}
	{A control costs 80,000 dollars per year and reduces expected annual loss from 120,000 to 70,000 dollars. What is the \textbf{BEST} conclusion?}
	{\item[\textbf{(A)}] It is justified because loss is reduced
		\item[\textbf{(B)}] It is justified because 120,000 is greater than 80,000
		\item[\textbf{(C)}] The control is not cost-justified based on loss reduction alone, since it saves 50,000 for an 80,000 cost
		\item[\textbf{(D)}] Cost is irrelevant to risk decisions}
	{C}
	{\item Compare reduction in expected loss to control cost.
		\item Non-financial factors may still apply, such as regulation.
		\item Decisions should be documented.}
	{Benefit must exceed cost, unless law or ethics say otherwise.}
	
	\hq{63}{Likelihood Inputs}
	{Which sources \textbf{BEST} support a likelihood estimate for a risk scenario?}
	{\item[\textbf{(A)}] Only the opinion of the loudest manager
		\item[\textbf{(B)}] Incident history, threat intelligence, control effectiveness, and expert judgment
		\item[\textbf{(C)}] A coin flip
		\item[\textbf{(D)}] Marketing material}
	{B}
	{\item Multiple sources improve reliability.
		\item Control strength affects likelihood.
		\item Document assumptions.}
	{Estimate with evidence, not volume.}
	
	\hq{64}{Business Impact Analysis Purpose}
	{What is the main purpose of a business impact analysis?}
	{\item[\textbf{(A)}] To rate employees
		\item[\textbf{(B)}] To identify critical processes and the impact of their disruption over time
		\item[\textbf{(C)}] To test network speed
		\item[\textbf{(D)}] To choose office furniture}
	{B}
	{\item BIA determines criticality and recovery priorities.
		\item It informs RTO and RPO.
		\item It uses business input, not only IT.}
	{BIA tells you what must come back first.}
	
	\hq{65}{BIA vs Risk Assessment}
	{How does a BIA differ from a risk assessment?}
	{\item[\textbf{(A)}] They are identical
		\item[\textbf{(B)}] A BIA considers only likelihood
		\item[\textbf{(C)}] A BIA focuses on the impact of disruption to processes, while a risk assessment considers threats, vulnerabilities, likelihood, and impact
		\item[\textbf{(D)}] A risk assessment ignores impact}
	{C}
	{\item BIA is impact-centered and process-focused.
		\item Risk assessment adds likelihood and causes.
		\item The two inform each other.}
	{BIA asks what hurts; risk assessment asks how likely.}
	
	\hq{66}{Maximum Tolerable Downtime}
	{The BIA shows that if a process is down more than 8 hours, the business faces severe consequences. What does this value represent?}
	{\item[\textbf{(A)}] Maximum tolerable downtime
		\item[\textbf{(B)}] Recovery point objective
		\item[\textbf{(C)}] Risk tolerance for phishing
		\item[\textbf{(D)}] Mean time between failures}
	{A}
	{\item MTD is the longest acceptable outage.
		\item RTO must be shorter than MTD.
		\item Dependencies must also be recoverable in time.}
	{RTO must fit inside MTD.}
	
	\hq{67}{Inherent Risk}
	{Which term describes risk before any controls or responses are considered?}
	{\item[\textbf{(A)}] Residual risk
		\item[\textbf{(B)}] Target risk only
		\item[\textbf{(C)}] Inherent risk
		\item[\textbf{(D)}] Accepted risk}
	{C}
	{\item Inherent risk is the raw exposure.
		\item Controls reduce it to residual risk.
		\item Comparing both shows control value.}
	{Inherent minus controls equals residual.}
	
	\hq{68}{Risk Register Quality}
	{Which risk register entry is \textbf{BEST} written?}
	{\item[\textbf{(A)}] Unpatched internet-facing server could be exploited, causing data loss and 2 days of outage; owner: operations director
		\item[\textbf{(B)}] Server issue
		\item[\textbf{(C)}] Bad things might happen
		\item[\textbf{(D)}] Fix patches}
	{A}
	{\item A clear risk statement links cause, event, and impact.
		\item An owner is named.
		\item Vague entries cannot be managed.}
	{Write the risk so a decision maker can act on it.}
	
	\hq{69}{Risk Scoring Consistency}
	{Different assessors rate the same risk very differently. What is the \textbf{BEST} fix?}
	{\item[\textbf{(A)}] Define and train on common rating criteria, then calibrate results across assessors
		\item[\textbf{(B)}] Average randomly
		\item[\textbf{(C)}] Use the highest rating always
		\item[\textbf{(D)}] Let each business unit invent its own scale}
	{A}
	{\item Shared scales produce comparable results.
		\item Calibration reduces bias.
		\item Consistency supports aggregation.}
	{Calibrate the ruler before you measure.}
	
	\hq{70}{Heat Map Limits}
	{What is a limitation of relying only on a heat map to report risk?}
	{\item[\textbf{(A)}] It uses colors
		\item[\textbf{(B)}] It can oversimplify and hide assumptions, so it should be supported with context and scenario detail
		\item[\textbf{(C)}] It is always inaccurate
		\item[\textbf{(D)}] It cannot be printed}
	{B}
	{\item Heat maps summarize but lose nuance.
		\item Boundary values can mislead.
		\item Narrative explains drivers and actions.}
	{Pictures need explanations.}
	
	\hq{71}{Risk Aggregation}
	{Many small risks are assessed in different departments. What is a benefit of aggregating them?}
	{\item[\textbf{(A)}] It hides risk owners
		\item[\textbf{(B)}] It reduces the number of controls needed
		\item[\textbf{(C)}] It reveals combined exposure and common causes that individual assessments may miss
		\item[\textbf{(D)}] It removes the need for appetite}
	{C}
	{\item Related risks can compound.
		\item Common causes suggest shared responses.
		\item Aggregation supports enterprise reporting.}
	{Small risks can add up to a big one.}
	
	\hq{72}{Using Assessment Results}
	{A risk assessment is complete. What should happen next?}
	{\item[\textbf{(A)}] File the report and take no action
		\item[\textbf{(B)}] Delete the underlying data
		\item[\textbf{(C)}] Wait one year before sharing
		\item[\textbf{(D)}] Communicate results to risk owners and decision makers and agree on responses}
	{D}
	{\item Assessments exist to support decisions.
		\item Owners decide on response.
		\item Timely communication keeps results relevant.}
	{An assessment nobody acts on has no value.}
	
	\hq{73}{Reassessment Triggers}
	{Which event is \textbf{MOST} likely to require a risk reassessment?}
	{\item[\textbf{(A)}] A new office plant
		\item[\textbf{(B)}] A holiday party
		\item[\textbf{(C)}] A change in parking rules
		\item[\textbf{(D)}] A significant change to a critical system or business process}
	{D}
	{\item Change alters threats, vulnerabilities, and controls.
		\item Reassess promptly when conditions change.
		\item Waiting for the annual cycle may miss new risk.}
	{Change triggers reassessment.}
	
	\hq{74}{Risk Velocity}
	{A risk could materialize within hours of a trigger, leaving very little time to respond. Which risk attribute does this describe?}
	{\item[\textbf{(A)}] Residual risk
		\item[\textbf{(B)}] Risk appetite
		\item[\textbf{(C)}] Velocity, meaning the speed at which the risk event unfolds
		\item[\textbf{(D)}] Inherent cost}
	{C}
	{\item Fast-moving risks need pre-planned responses.
		\item Velocity complements likelihood and impact.
		\item Monitoring and automation help shorten reaction time.}
	{The faster the risk, the earlier you must prepare.}
	
	% ============================================================
	\section{Incidents, Issues and the Risk Lifecycle}
	% ============================================================
	
	\hq{75}{Event vs Incident}
	{Which statement \textbf{BEST} distinguishes an event from an incident?}
	{\item[\textbf{(A)}] They mean the same thing
		\item[\textbf{(B)}] An incident is harmless
		\item[\textbf{(C)}] An event is any observable occurrence, while an incident is an event that causes or could cause harm or disruption
		\item[\textbf{(D)}] An event always means a breach}
	{C}
	{\item Many events are routine.
		\item Incidents require response.
		\item Breaches are incidents with confirmed unauthorized access.}
	{Events happen constantly; incidents hurt.}
	
	\hq{76}{First Incident Priority}
	{A serious incident is confirmed. What should be the \textbf{FIRST} priority?}
	{\item[\textbf{(A)}] Find someone to blame
		\item[\textbf{(B)}] Rebuild all systems immediately without analysis
		\item[\textbf{(C)}] Issue a press release before facts are known
		\item[\textbf{(D)}] Contain the incident to limit damage while preserving evidence}
	{D}
	{\item Containment stops the spread.
		\item Evidence supports investigation and legal needs.
		\item Recovery follows containment.}
	{Contain, then investigate, then recover.}
	
	\hq{77}{Incident Response Plan}
	{What is the \textbf{MOST} important attribute of an effective incident response plan?}
	{\item[\textbf{(A)}] A very long document
		\item[\textbf{(B)}] Defined roles, escalation paths, and procedures that have been tested
		\item[\textbf{(C)}] Secrecy from staff
		\item[\textbf{(D)}] Use only by IT}
	{B}
	{\item Clear roles prevent confusion.
		\item Testing reveals gaps.
		\item Plans should involve business and legal.}
	{A plan untested is a plan unproven.}
	
	\hq{78}{Lessons Learned Timing}
	{When should a lessons-learned review be held after an incident?}
	{\item[\textbf{(A)}] Never
		\item[\textbf{(B)}] Soon after recovery, while details are fresh, with findings fed into the risk register
		\item[\textbf{(C)}] After one year
		\item[\textbf{(D)}] Only if there was a fine}
	{B}
	{\item Fresh recall improves quality.
		\item Findings drive improvement.
		\item Update scenarios and controls.}
	{Learn while the incident is still vivid.}
	
	\hq{79}{Issue Management}
	{An audit identifies a control weakness. How should it be handled in the risk process?}
	{\item[\textbf{(A)}] Ignore it until the next audit
		\item[\textbf{(B)}] Close it without action
		\item[\textbf{(C)}] Share it only with auditors
		\item[\textbf{(D)}] Record it as an issue, assess its risk, assign an owner, and track the response to closure}
	{D}
	{\item Issues are inputs to risk assessment.
		\item Owners and deadlines enable follow-up.
		\item Track until verified closed.}
	{Every issue needs an owner and an end date.}
	
	\hq{80}{Anomaly Response}
	{A monitoring system shows unusual outbound traffic at night. What should be done \textbf{FIRST}?}
	{\item[\textbf{(A)}] Dismiss it as noise
		\item[\textbf{(B)}] Disconnect the entire company
		\item[\textbf{(C)}] Investigate to determine whether it is a true incident, following the response procedure
		\item[\textbf{(D)}] Wait for a customer complaint}
	{C}
	{\item Triage determines severity.
		\item Procedures guide proportionate response.
		\item Ignoring anomalies can allow breaches to grow.}
	{Investigate before you panic or dismiss.}
	
	\hq{81}{Risk Management Lifecycle}
	{Which sequence \textbf{BEST} reflects the IT risk management lifecycle?}
	{\item[\textbf{(A)}] Report, ignore, forget, repeat
		\item[\textbf{(B)}] Assess once and file
		\item[\textbf{(C)}] Respond, then identify
		\item[\textbf{(D)}] Identify, assess, respond, monitor, and report, repeating continuously}
	{D}
	{\item Risk management is iterative.
		\item Monitoring feeds new identification.
		\item Reporting supports governance.}
	{The cycle never ends.}
	
	\hq{82}{Near Miss}
	{A phishing email almost led to credential theft but was caught by an alert employee. What should be done?}
	{\item[\textbf{(A)}] Ignore it because nothing was lost
		\item[\textbf{(B)}] Punish the employee who received it
		\item[\textbf{(C)}] Keep it secret
		\item[\textbf{(D)}] Treat it as a near miss, learn from it, and update controls and awareness}
	{D}
	{\item Near misses reveal weaknesses cheaply.
		\item Recognize good reporting behavior.
		\item Update indicators and training.}
	{Near misses are free lessons.}
	
	\hq{83}{Third-Party Incident Ownership}
	{Your outsourced service provider suffers an outage affecting your customers. Who is \textbf{ULTIMATELY} accountable to your customers?}
	{\item[\textbf{(A)}] Your organization
		\item[\textbf{(B)}] The service provider alone
		\item[\textbf{(C)}] The customers
		\item[\textbf{(D)}] No one}
	{A}
	{\item Accountability to customers remains with you.
		\item Contracts can allocate cost but not accountability.
		\item Oversight must match criticality.}
	{Customers hold you accountable, not your vendor.}
	
	\hq{84}{Regulatory Reporting Incident}
	{An incident may trigger regulatory reporting. What is the \textbf{BEST} practice?}
	{\item[\textbf{(A)}] Pre-define legal review and reporting procedures so decisions are fast and accurate
		\item[\textbf{(B)}] Decide after press coverage
		\item[\textbf{(C)}] Avoid involving legal
		\item[\textbf{(D)}] Let each employee decide}
	{A}
	{\item Deadlines may be very short.
		\item Prepared playbooks reduce errors.
		\item Consistent messages protect the organization.}
	{Prepare the notification process before you need it.}
	
	\hq{85}{Root Cause}
	{After repeated outages, the team fixes each symptom but they recur. What is the \textbf{BEST} next step?}
	{\item[\textbf{(A)}] Add more staff to the help desk
		\item[\textbf{(B)}] Perform root cause analysis and address the underlying cause
		\item[\textbf{(C)}] Restart systems more often
		\item[\textbf{(D)}] Accept outages as normal}
	{B}
	{\item Symptoms are not causes.
		\item Root cause fixes prevent recurrence.
		\item Update the risk register with findings.}
	{Fix the cause, not the symptom.}
	
	\hq{86}{Risk Monitoring Ownership}
	{Who should \textbf{PRIMARILY} monitor day-to-day performance of a risk response within a department?}
	{\item[\textbf{(A)}] The risk owner and control owners, with oversight from the second line
		\item[\textbf{(B)}] The external auditor only
		\item[\textbf{(C)}] The board only
		\item[\textbf{(D)}] No one}
	{A}
	{\item First line operates and watches controls.
		\item Second line challenges and aggregates.
		\item Third line assures independently.}
	{Monitoring starts where the work happens.}
	
	% ============================================================
	\section{Mixed Exam-Style Scenarios}
	% ============================================================
	
	\hq{87}{Best Next Step}
	{A new cloud service is already live and used by finance, but was never assessed. What should the risk practitioner do \textbf{FIRST}?}
	{\item[\textbf{(A)}] Understand what data and processes are involved and assess the risk with the business owner
		\item[\textbf{(B)}] Shut it down immediately
		\item[\textbf{(C)}] Ignore it because it is already live
		\item[\textbf{(D)}] Report the finance team to HR}
	{A}
	{\item Understand before acting.
		\item The owner provides business context.
		\item Blocking may be considered after assessment.}
	{Assess, then decide; do not react blindly.}
	
	\hq{88}{Most Important Factor}
	{Which is the \textbf{MOST} important factor when determining a risk's priority?}
	{\item[\textbf{(A)}] The technology involved
		\item[\textbf{(B)}] How interesting the risk is
		\item[\textbf{(C)}] The name of the person who found it
		\item[\textbf{(D)}] The potential impact on business objectives}
	{D}
	{\item Priority follows business impact.
		\item Technical detail supports but does not decide priority.
		\item Objectives define what matters.}
	{Business objectives set the priority.}
	
	\hq{89}{Stakeholder Disagreement}
	{The business owner disagrees with the risk rating. What should the risk practitioner do?}
	{\item[\textbf{(A)}] Change the rating immediately
		\item[\textbf{(B)}] Refuse to discuss
		\item[\textbf{(C)}] Review the assumptions and evidence together, revise if justified, and escalate unresolved disagreement
		\item[\textbf{(D)}] Delete the risk}
	{C}
	{\item Dialogue may reveal new information.
		\item Ratings change only with evidence.
		\item Unresolved disputes go to governance.}
	{Discuss the evidence, not the preference.}
	
	\hq{90}{Budget Cut}
	{Management cuts the security budget. What is the risk practitioner's \textbf{BEST} action?}
	{\item[\textbf{(A)}] Assess the impact on risk exposure and present it to decision makers
		\item[\textbf{(B)}] Quietly accept the cut
		\item[\textbf{(C)}] Resign
		\item[\textbf{(D)}] Hide the impact}
	{A}
	{\item Decision makers need to understand consequences.
		\item Impact is expressed in business terms.
		\item They may accept the risk knowingly.}
	{Show what the cut costs in risk.}
	
	\hq{91}{Audit Disagreement}
	{An auditor rates a control as ineffective, but the control owner disagrees. What should happen?}
	{\item[\textbf{(A)}] Ignore the audit
		\item[\textbf{(B)}] Remove the control
		\item[\textbf{(C)}] Review the evidence together and resolve factual differences, with risk assessed on the final finding
		\item[\textbf{(D)}] Threaten the auditor}
	{C}
	{\item Facts should be agreed first.
		\item Independence must be respected.
		\item Residual risk should be updated.}
	{Evidence settles disputes.}
	
	\hq{92}{Automation Risk}
	{A robotic process automation tool will handle payments. Which risk consideration is \textbf{MOST} important?}
	{\item[\textbf{(A)}] The robot's color
		\item[\textbf{(B)}] The number of staff that can be laid off
		\item[\textbf{(C)}] Marketing claims
		\item[\textbf{(D)}] Access controls, change control, monitoring, and clear ownership for the automated process}
	{D}
	{\item Automation can scale errors quickly.
		\item Bots need identities and least privilege.
		\item Ownership and monitoring remain essential.}
	{Automation still needs controls and an owner.}
	
	\hq{93}{Merger Integration}
	{After a merger, two networks will be connected. What should the risk practitioner recommend \textbf{FIRST}?}
	{\item[\textbf{(A)}] Assess the risk posture of the acquired environment before connecting
		\item[\textbf{(B)}] Connect immediately to save time
		\item[\textbf{(C)}] Ignore the acquired controls
		\item[\textbf{(D)}] Merge all accounts without review}
	{A}
	{\item Unknown weaknesses can spread.
		\item Assessment informs integration plans.
		\item Segmentation can limit exposure.}
	{Do not connect what you have not assessed.}
	
	\hq{94}{Remote Work Shift}
	{The company moves to permanent remote work. What should the risk practitioner do?}
	{\item[\textbf{(A)}] Reassess risks, update policies and controls for home environments, and communicate changes
		\item[\textbf{(B)}] Assume the previous risk profile still applies
		\item[\textbf{(C)}] Prohibit all home networks
		\item[\textbf{(D)}] Wait for incidents}
	{A}
	{\item Work location changes threats and controls.
		\item Policies and training must adapt.
		\item Reassessment supports good decisions.}
	{New ways of working need new risk views.}
	
	\hq{95}{Metrics for Management}
	{A manager asks how the risk program is doing. Which information is \textbf{BEST}?}
	{\item[\textbf{(A)}] The number of meetings held
		\item[\textbf{(B)}] Total pages in the risk policy
		\item[\textbf{(C)}] Trends in key risk indicators against appetite, progress on response plans, and notable changes
		\item[\textbf{(D)}] The size of the risk team}
	{C}
	{\item Outcome measures show progress.
		\item Comparison to appetite gives context.
		\item Activity counts do not show value.}
	{Report outcomes, not activity.}
	
	\hq{96}{Whistleblower Issue}
	{An employee anonymously reports control override by a senior executive. What should be done?}
	{\item[\textbf{(A)}] Escalate through governance channels such as the audit committee, preserving confidentiality and evidence
		\item[\textbf{(B)}] Tell the executive who reported
		\item[\textbf{(C)}] Ignore the report
		\item[\textbf{(D)}] Delete the evidence}
	{A}
	{\item Independent oversight is needed.
		\item Protect reporters from retaliation.
		\item Evidence must be preserved.}
	{Serious allegations go to independent oversight.}
	
	\hq{97}{Cost of Control}
	{A proposed control is expensive and reduces a low risk slightly. What is the \textbf{BEST} recommendation?}
	{\item[\textbf{(A)}] Implement it because security always comes first
		\item[\textbf{(B)}] Reject it without explanation
		\item[\textbf{(C)}] Consider alternative, lower-cost responses or accept the risk with owner approval
		\item[\textbf{(D)}] Hide the decision}
	{C}
	{\item Controls must be proportionate.
		\item Alternatives may meet the objective.
		\item Acceptance should be documented.}
	{Proportionality beats perfection.}
	
	\hq{98}{New Technology Adoption}
	{The business wants to adopt an emerging technology quickly. What should the risk practitioner do?}
	{\item[\textbf{(A)}] Block it until it is mature
		\item[\textbf{(B)}] Engage early to identify risks and recommend proportionate controls without blocking value
		\item[\textbf{(C)}] Ignore it
		\item[\textbf{(D)}] Approve without review}
	{B}
	{\item Early engagement enables safe adoption.
		\item Controls should fit the risk.
		\item Risk practitioners enable the business.}
	{Enable the business safely.}
	
	\hq{99}{Ultimate Accountability}
	{Which statement about accountability for risk is \textbf{MOST} accurate?}
	{\item[\textbf{(A)}] It can be fully transferred to vendors
		\item[\textbf{(B)}] Accountability for risk stays with the organization and its owners even when activities are outsourced
		\item[\textbf{(C)}] It belongs to the risk practitioner alone
		\item[\textbf{(D)}] It disappears after insurance}
	{B}
	{\item Contracts transfer tasks, not accountability.
		\item Owners must oversee outsourced work.
		\item Insurance covers cost, not responsibility.}
	{You can delegate tasks, not accountability.}
	
	\hq{100}{Exam Mindset}
	{On the exam, several answers seem technically correct. Which approach usually leads to the \textbf{BEST} answer?}
	{\item[\textbf{(A)}] Choose the option that best aligns with business objectives, ownership, and a risk-based, proportionate approach
		\item[\textbf{(B)}] Choose the most technical option
		\item[\textbf{(C)}] Choose the most expensive option
		\item[\textbf{(D)}] Choose the option that acts fastest regardless of context}
	{A}
	{\item ISACA favors business alignment and risk-based decisions.
		\item Understand before acting.
		\item Owners decide; practitioners advise.}
	{Think like an advisor to the business, not a technician.}
	
	% ============================================================
	\section{Conclusion}
	% ============================================================
	
	\begin{frame}{Exam Mindset Summary}
		\begin{block}{Remember These Core Principles}
			\begin{itemize}
				\item \textbf{Business First:} Risk is defined by its effect on business objectives.
				\item \textbf{Owners Decide:} Practitioners inform; owners accept, mitigate, transfer, or avoid.
				\item \textbf{Three Lines:} The first line owns risk, the second oversees, and the third gives independent assurance.
				\item \textbf{Ethics Matter:} Honesty, independence, and confidentiality underpin credible risk advice.
				\item \textbf{Evidence Over Opinion:} Use data, scenarios, and calibrated ratings.
				\item \textbf{Proportionality:} Controls should fit the risk and its cost.
				\item \textbf{Accountability Cannot be Outsourced:} Vendors and cloud providers share duties, but accountability remains with the organization.
			\end{itemize}
		\end{block}
	\end{frame}
	
	\begin{frame}{Thank You}
		\begin{center}
			{\Huge \textbf{Good Luck on Your CRISC Exam!}}
			
			\vspace{1cm}
			
			\normalsize
			Real exam questions test your ability to apply these concepts in practical, real-world scenarios.
		\end{center}
	\end{frame}
	
\end{document}